\chapter{C1.1: an opening-uniform full normal and a new angular circuit}
\label{component:p9afbe73899af}
\noindent\textbf{Component }\path{p9afbe73899af}\par
{\small\noindent Source: \path{R13/advances/c1_small/small_u_harmonic_reduction.tex}\par}
{\small\noindent SHA256: \path{9afbe73899afc5182fd991d7b285b98dae466d7ecca55667cf04a3e5e251a889}\par}
{\small\noindent Inventory use: SC1--SC9; exact source cancellation, full-parent reduction and remote/state proof | SC2-SC20 | C1\_\allowbreak{}BARE\_\allowbreak{}OTHER\_\allowbreak{}TERMINALS; precise subordinate uses retained below\par}
\medskip
\section{C1.1: an opening-uniform full normal and a new angular circuit}
\label{p9afbe73899af:sec:r7-c1-small}
This section concerns the original C1.1 source, not an enlarged family.
All coefficients are divided by $d>0$. Put
\[
 b=\eta^{-1/2},\quad g=\sqrt\zeta,\quad q=k^2\eta,
 \qquad x=b^2q,\quad a=b^2(1-r).
 \tag{SC1}
\]
Here $b,g$ are positive infinitesimals, $q>0$ is bounded, and
$\overline q\le1/2$. The new bounded two-normal reduction below is
asserted for $q$ a positive unit. The exact polynomial identities
remain valid more generally, but their indicated boundedness argument
would not control $g/\sqrt q$ when $q$ is infinitesimal.
No all-tests closure is inferred from the reduction.

Keep $K,R,T,M,J,E,Q,P_4,P_d,P_\alpha$ from (NC1)--(NC2), and set
\[
\begin{aligned}
 c_b&=b^2C/d=\frac{K(1+b^2r)E}{2(1-r)P_d},&
 a_b&=b^2A/d=\frac{P_4}{2rP_d},\\
 D_b&=D/(db^2),&\alpha_b&=\alpha/(db^2g^2),\\
 h_2&=\frac{Kq(1+b^2r)(1-tg^2)(1+qg^2)}{2P_d},&
 h_b&=\frac{Kq(1+b^2r)[1+(1-t)g^2]}{2P_d}.
\end{aligned}
\]
Thus $m/d=h_2/(g^2M)$ and $\beta/d+1/2=h_b/(g^2t)$.
The exact relation
\[
 h_b-h_2=\frac{Kq(1+b^2r)E g^2}{2P_d}
 \tag{SC2}
\]
removes the apparent opening dependence from the first normal of the
paired $m,\beta$ currents. The quantities $D_b,\alpha_b,h_2,h_b$ are
bounded; $\alpha_b,h_2,h_b$ are units in the stated $q$-unit domain.
Their residues are respectively $(1-2q)/2,q^2/2,q/2,q/2$, with
bounded source letters on the right understood at residue.

\paragraph{Exact full-parent reduction.}
For bounded affine tests define
\[
 Q_b=(W+uY)/b,\quad Y_b=Y/b^2,\quad p=t/R,\quad
 N=\frac{(X+uZ)/b-pvQ_b}{g},
\]
\[
 H=uQ_b/b,\quad L=vQ_b/g,\quad V=uvZ/b,\quad
 B_b=gN+pvQ_b,\quad D_u=1+u^2,\quad D_v=1+v^2.
\]
All roots use only pure-source inverses. Retain the two nonnegative
original packets
\[
 \frac{4h_2t}{MR}N^2,
 \qquad \frac{4h_bM}{tR}D_vN^2.
\]
The complete original parent divided by $db^2$, minus these packets,
is exactly
\begin{align*}
 G={}&\frac4R\{(h_2+h_bD_v)N^2+2(h_bD_v-h_2)LN
              +[h_2(1-p)+h_bpD_v]L^2\}\\
 &+4D_v\{c_b(H-D_uY_b)^2+a_bD_uY_b^2\}
   +4D_bJ_b^2+4\alpha_b g^2D_u(vb^2Y_b-Z)^2+2K_b,
                                                        \tag{SC3}\\
 J_b={}&b^2(H-D_uY_b)+ubvB_b-D_uvZ,\\
 K_b={}&b^2(H-D_uY_b)^2+2buv(H-D_uY_b)B_b+u^2v^2B_b^2\\
 &+D_uv^2b^2Y_b^2+D_uV^2
       -2D_uV(Q_b+vB_b)+2D_uu^2vZY_b.
\end{align*}
These are entire quadratic polynomials in the state; no restriction
of a state normal has been made. In particular, (SC3) is a direct
input for collecting both complete first normals and their full hard
quadratic, rather than an identity solely on a contact graph.

\paragraph{A harmonic normal removes both $r$ openings.}
Put
\[
 s=c_b+a_b,\quad \rho=c_b/s,\quad h=c_ba_b/s,
 \quad \rho_0=Kr/R,\qquad
 N_y=\sqrt{s}\,(D_uY_b-\rho_0H).
\]
Then the complete $C,A$ line in (SC3), divided by $4D_v$, equals
\[
 \left[N_y+\sqrt{s}(\rho_0-\rho)H\right]^2+hH^2
 -\frac{a_bu^2}{D_u}
           \left(\frac{N_y}{\sqrt{s}}+\rho_0H\right)^2.
 \tag{SC4}
\]
The inverse $1/D_u$ is pure testing, and every coefficient in (SC4)
is bounded uniformly through either $r\to0$ or $r\to1$.
Indeed $s$ is bounded below, $h$ is a positive unit, and
$a_b/s$, $a_b\rho_0/\sqrt{s}$ and $a_b\rho_0^2$ are bounded.
There is no mixed harmonic denominator.

Here is the exact source cancellation that also centers the normal
without an opening loss. Define
\[
 S_0=Kr(1+b^2r)E+(1-r)P_4,
\]
\[
 T_0=q^2(1-tg^2)^2+qrg^2+2qtg^2-qg^2-2q+r.
\]
Then
\[
 \rho=\frac{Kr(1+b^2r)E}{S_0},\qquad
 Kr-\rho R=-\frac{b^2r(1-r)K T_0}{S_0},\qquad
 h=\frac{K(1+b^2r)E P_4}{2P_dS_0}.             \tag{SC5}
\]
Both $S_0$ and the other denominators are positive units. Also
$s\asymp[r(1-r)]^{-1}$, so $\sqrt{s}(\rho_0-\rho)$ is
infinitesimal without a lower bound on either opening.
At the actual contact,
\[
 (u,v)=(b\sqrt q,\sqrt q g),\quad
 Q_b=RgZ,\quad Y_b=r\sqrt q gZ,
\]
and both $N,N_y$ vanish exactly. Thus (SC3)--(SC5) remove every
unbounded source weight while retaining the actual contact and every
state normal.

\paragraph{The leading angular circuit.}
For $u,v$ infinitesimal, the leading scalar after the harmonic state
normal has been removed is twice
\[
 (1-q)H^2+qL^2+V^2+(1-2q)(vZ)^2+q^2(gZ)^2-2VQ_b.
\]
Here the leading relation $H(vZ)=VQ_b$ is exact for the displayed
roots. It splits as
\[
 (1-2q)(H-vZ)^2+
 qH^2+qL^2+V^2+q^2(gZ)^2-4qVQ_b.              \tag{SC6}
\]
The final four positive terms form one AM--GM circuit, since
\[
 (\sqrt qH)(\sqrt qL)V(qgZ)=(qVQ_b)^2.
\]
On the positive negative-core chart its radical factors separate
between the source and testing fields. This proves positivity of the
leading scalar, including $q=1/2$, where its separate binomial drops
out. The leading hard normal coefficients are $4q$ on $N$ and $4$
on $N_y$. This organizing form alone is not yet a certificate for
all the higher terms in (SC3).

A useful exact conductor can be stated without comparing any source
and testing rate. Normalize $Z=1$ on its pure-testing state chart,
put $c=\sqrt q$, $w=(W+uY)/R$, and now use the contact-exact roots
\[
 H_0=uw/b^2,\ L_0=vw/(bg),\ V_0=uv/b,\ D_0=c^2g,
 \quad A_0=cH_0,\ B_0=cL_0,\ C_0=V_0.
\]
Their contact is $(u,v,w)=(bc,cg,bg)$. Set
\[
 E_0=H_0-v,\quad
 S_v=A_0-D_0-cE_0,\quad
 S_u=C_0-D_0-S_v,\quad S_w=B_0-D_0-S_v.
\]
Direct multiplication gives
\[
\begin{aligned}
 H_0(u-bc)&=(u-bc)E_0+bS_u,\\
 H_0(v-cg)&=(v-cg)E_0+(v/c)S_v,\\
 H_0(w-bg)&=(w-bg)E_0+(bg/c)S_w.
\end{aligned}                                                     \tag{SC7}
\]
Only the positive-unit inverse $1/c$ occurs. These identities lift
contact-vanishing coefficients times $H_0$ to the displayed root
module with bounded coefficients on the small pure-testing chart.
They do not alone prove a full scalar square lift: a first-jet
syzygy among the redundant angular roots still has to be allocated,
particularly when the coefficient of $E_0^2$ vanishes at $q=1/2$.

\paragraph{Exact small-cost angular pivots.}
The four-term circuit has the exact symmetric decomposition
\[
 \sum_{i=1}^4A_i^2-4\sqrt{A_1A_2A_3A_4}
 =\frac13\sum_{i<j}(A_i-A_j)^2
  +\frac23\sum_{\{ij\}\mid\{kl\}}
       (\sqrt{A_iA_j}-\sqrt{A_kA_l})^2,
\]
where the second sum is over the three pair partitions. Use
$(A_1,A_2,A_3,A_4)=(A_0,B_0,C_0,D_0)$ above and assume the positive
chart $u,v,w>0$. Then the following identities are exact:
\[
\begin{aligned}
 c\sqrt{uw}\,E_0
 &=\sqrt{uw}(A_0-D_0)
      +b\sqrt{cg}(\sqrt{A_0D_0}-\sqrt{B_0C_0}),\\
 \sqrt{vw}(u-bc)
 &=\frac{b^{3/2}}{\sqrt c}
             (\sqrt{A_0C_0}-\sqrt{B_0D_0}),\\
 \sqrt{uw}(v-cg)
 &=b\sqrt{g/c}
             (\sqrt{B_0C_0}-\sqrt{A_0D_0}),\\
 \sqrt{uv}(w-bg)
 &=\frac{b^{3/2}\sqrt g}{c}
             (\sqrt{A_0B_0}-\sqrt{C_0D_0}).
\end{aligned}                                                     \tag{SC9}
\]
The coefficient factors on the right are infinitesimal; all radicals
separate between the source and testing fields. Consequently any fixed
positive real multiples of the squares of the four left-hand roots can
be inserted as additional positive Gram pivots at an infinitesimal cost
in the displayed strict circuit Gram. To see this without a limiting
assertion, write the four identities as $P=TZ$, with $Z$ the displayed
circuit roots and every entry of $T$ in the bounded maximal ideal. Then
$Z^TGZ=\tau P^TP+Z^T(G-\tau T^TT)Z$, and the second matrix retains
the strictly positive real residue of $G$ for every fixed real $\tau>0$.
This is a genuine full-rates Gram mechanism. It still requires showing
that the entire first-normal and scalar remainders of (SC3) allocate on
these weighted roots; the identities do not make that last membership
automatic.

\paragraph{A quantified remote improvement when $e$ is unbounded.}
The source identity $e^2=q/g^2$ and (NI3) imply that $m,\beta$ are
both unbounded when $e$ is unbounded. Three inverse-cycle rotations
send this adjacent pair to $C,A$ and transform tests as
$(u,v)\mapsto(v,1/u)$. Therefore the adjacent-event theorem covers
all tests with $v$ not infinitesimal. Combined with the previous
$C,A$ admission, only $u,v$ infinitesimal remain in this $e$ type.
This includes every $q$-unit source considered above.

There is also a complete state-chart admission there. Normalize the
pure-testing state by a coordinate of maximal absolute value. If
$W$ is a unit, cap $C,A,m,\beta$ at a sufficiently large positive
real $\mu$, retain all excesses, and leave every other event as its
original nonnegative square. At $u=v=X=Y=0$, the quadratic part in
$(u,v,X,Y)$ is
\[
 2(uW-Y)^2+2(X+uZ)^2-4uvWZ+
 4\mu\{(uW-Y)^2+Y^2+(vW-X-uZ)^2+(X+uZ)^2\}.
 \tag{SC8}
\]
The event part is strictly positive for every real $W\ne0$ and bounded
real $Z$: its zero equations give successively $Y=u=X=v=0$. A sufficiently
large finite $\mu$ therefore absorbs the entire signed quadratic core.
The exact remainder has degree at least three in these four small
roots, with bounded coefficients in $W,Z$, so it allocates with zero
residue. This proves the entire original parent on this chart.
If the core has negative sign and $W$ is infinitesimal, then $Z$ is
a unit, since otherwise the core residue forces $X=Y=0$, contrary
to maximal-coordinate normalization. Thus the remaining state chart
has $Z=1$, $W$ infinitesimal, and both testing half-angles infinitesimal.
This is a full quantified narrowing, not closure of C1.1.

\paragraph{Completion and exact outstanding obligation.}
The following section supplies a finite scalar square lift and both
complete first normals for the $q$-unit C1.1 subdomain, including
$D=0$. It uses a contact-centered $D$ packet followed by the geometric
mean conductor; an unweighted $E_0^2$ perturbation by itself is not
controlled by the circuit uniformly at arbitrary relative rates.
For $q$ infinitesimal, the uniform normal and protected contact lift
in Sections~\ref{pc376284fc7ae:sec:r7-c1-q-inf-local}--\ref{pd2ecc6ac1891:sec:r7-c1-q-inf-remote}
resolve the additional $g/\sqrt q$ issue and all remote tests.
No assumption of bounded $e$ or $1/e$ is imposed. The original
checker verifies thirty-nine whole-polynomial and source identities;
the later modules include their own exact checks and scoped review.
